问题分析围绕共同的资源表示展开。原始计划按频段、首次时间、间隔和次数展开为完整重复事件，再映射到时间—频率资源格；计划层识别冲突，事件层记录重叠次数，资源格层写出互斥约束。三层表示相互对应：Q1 的冲突边传给 Q2，Q2 的验证方案成为 Q3 的冻结背景，所有结果回溯同一份预处理结果。

\CUMCMAnalysisLead{一}
关键是把计划记录还原为可比较的事件集合。对第 $i$ 条计划，依据首次使用时刻、单次时长、间隔和次数递推全部事件，统一采用半开时间区间。先用频段区间相交筛选计划对，再检查候选对中是否存在时间相交的重复事件；计划层只保留一条无向冲突边，事件级重叠次数另行保存。随后把事件占用转换为离散资源格，用倒排关系复核冲突清单。150 条计划展开 1\,300 个事件，枚举 11\,175 个候选对，得到 297 对冲突装备和 431 次事件级重叠，其中 B--C 类为 181 对。输出的冲突图同时作为问题一结果和问题二硬约束输入。

\CUMCMAnalysisLead{二}
难点是把“少撤销、保护高等级、少调整”变成可审计的优先级。对每条计划生成保持、频移、时移和撤销候选；候选均重新展开事件并形成资源格集合，越界者先剔除。一条计划只能选一种动作，一个资源格至多由一个候选使用，于是得到 0--1 资源格模型。目标不采用未经授权的加权和，而按 $(R,R_A,R_B,M,M_A,M_B,S_{10})$ 逐层固定；每层记录可行值、下界或阈值不可行证据，只有上下界闭合才写成已证明结论。4\,582 个合法候选动作对应 52\,939 条资源格团约束，前六层闭合为 $(6,0,4,126,16,34)$；$S_{10}=775$ 仍只是该前缀下的可行上界。连续区间与离散资源格回读均零冲突，144 个活动计划及 15\,816 个占用格成为问题三的冻结输入。

\CUMCMAnalysisLead{三}
研究固定问题二方案后的剩余容量，不允许新增计划反向改变 Q2 布局。继承 Q2 的 144 个活动计划及其资源格占用，按 C 类共同结构枚举合法的时间起点和频段起点；每个起点均展开完整事件，并通过资源域和固定背景冲突检查。候选按逐格占用筛选，而不是用剩余面积估计；随后用资源格集合装填和候选两两冲突两种编码交叉求解。52\,136 个完整 C 类起点中有 2\,334 个与冻结背景不冲突，两种模型均得到 141 台新增计划，且“至少 142 台”不可行。因此，141 只是在当前 Q2、资源域和 C 类模板下的最大值。

\CUMCMAnalysisLead{四}
问题四的变化不在于重新定义冲突，而在于放开 C 类计划的重复间隔。间隔一旦改变，第二次及以后的所有事件都会联动移动，因此不能只修补某一个事件，也不能用平均占用量代替完整展开。处理方式是从问题一的 150 条原始计划独立起算，枚举保持、频移、时移、撤销以及 C 类间隔调整动作；对每个动作递推全部事件，先剔除越界候选，再用资源格互斥约束统一判定。这样可把“允许更灵活的周期结构”转化为有限的整数动作集合。基准域共生成 6\,103 个合法候选动作，支配替换后保留 5\,533 个活动候选。3 台撤销的可行见证与 2 台及以下六种类别分支的不可行证据分别形成上下界，闭合当前有限模型下的最小撤销层。需要特别区分的是，$R=3$ 之外的类别撤销、调整数量和末级变化量仍只是当前条件候选，不能因求解器返回可行解就改写为完整字典序最优。

四问接口由预处理和边界验证共同固定：原始字段只在预处理阶段解释一次，重复事件采用同一半开区间，资源格使用同一时间域和频段域。Q2 的 $H=638,643,648$ 检验均给出最小撤销数 6；Q4 在三个边界情景均找到通过验证的三撤销候选。由此，Q1 的冲突事实、Q2 的条件字典序方案、Q3 的固定背景最大值和 Q4 的最小撤销层形成可追踪证据链，各问仍按照各自的输入接口和证明范围报告。
